\documentclass[11pt]{article}
\usepackage[a4paper,margin=27mm]{geometry}
\usepackage[T1]{fontenc}
\usepackage{lmodern,microtype,amsmath,amssymb,amsthm,mathtools,enumitem}
\usepackage[hidelinks,hypertexnames=false]{hyperref}
\newtheorem{theorem}{Theorem}
\newtheorem{lemma}[theorem]{Lemma}
\newtheorem{proposition}[theorem]{Proposition}
\newtheorem{corollary}[theorem]{Corollary}
\theoremstyle{definition}\newtheorem{definition}[theorem]{Definition}
\theoremstyle{remark}\newtheorem{remark}[theorem]{Remark}
\newcommand{\R}{\mathbb R}
\newcommand{\Dscr}{\mathcal D}
\newcommand{\supp}{\operatorname{supp}}
\newcommand{\debt}{\mathrm{debt}}
\newcommand{\cred}{\mathrm{credit}}
\newcommand{\sil}{\mathrm{silence}}
\newcommand{\Bfam}{\mathcal B}
\newcommand{\Jid}{\mathcal J}
\newcommand{\Iid}{\mathcal I}
\newcommand{\Uf}{\mathcal U}
\newcommand{\UU}{\mathfrak U}
\newcommand{\Alg}{\mathfrak A}
\newcommand{\Hi}[3]{\mathrm H^{#1}_{#2}(#3)}
\newcommand{\pw}{\mathrm{pw}}
\title{Three Complete Rankings Cannot Cover}
\author{}\date{11 October 2026}
\begin{document}\maketitle

\noindent
This note continues FSE (\emph{Fineness Spreads Along Edges}) and CCF
(\emph{Closed Choice Forces Fineness}), and uses, through them, results of
CBOI (\emph{Closed Choice Beyond Ordinary Integrals}). For one consequence it
also uses TAW (\emph{Two Additive Worlds Without Completeness}). It solves
problem~P1 of FSE, Section~1.7: with three worlds, rankings that are
complete preorders satisfying (G) and (IB), and the selection axiom (S),
the vacuous model $V$ is dominated.

\section{Summary and consequences}\label{sec:summary}

\begin{enumerate}[leftmargin=*,itemsep=3pt]
\item \emph{Main theorem} (Theorem~\ref{thm:main}). Let $n=3$, let every
ranking be a complete preorder satisfying (G) and (IB), and assume (S). Then
$V$ is dominated, and hence dominated by a selected coherent model. This
answers P1.
\item \emph{Full coverage needs incomparabilities}
(Corollary~\ref{cor:coverage}). Let $n=3$, and assume (G) and (IB) at every
world and (S). If every coherent model is undominated, or merely $V$ is,
then the ranking of at least one world is incomplete. By TAW, moreover, at
most one world has a one-sided weak extension of CBOI's finitely additive
form.
\item \emph{No saturated cut} (Theorem~\ref{thm:cut}). Let world $v$ be
complete, and assume (G) at worlds $v$ and $w$ and (S). For every family of
measurable subsets of $D_{vw}$ that contains a non-null subset of every band
of the edge $[e_v,e_w]$, $D_w$ is reached from $D_v$ by a finite chain of descents in
$\preceq_v$ and of perturbations on members of the family. Neither (IB) nor
any hypothesis on $V$ is used.
\item \emph{Counting lemma} (Lemma~\ref{lem:count}). Let $\preceq_v$ be a
preorder with weak (G) and (IB). If such a chain has length $\ell\ge2$ and
all its perturbation sets lie in $F$, then there are no $\ell$ pairwise disjoint
subsets of $D_{vw}\setminus F$ each of which carries a full change from
credit to debt that outweighs the full change on every perturbation set.
Neither completeness nor (S) is used.
\item \emph{Rich far ends} (Proposition~\ref{prop:rich}). Under the
hypotheses of P1, if $V$ is undominated, the point of $[e_1,e_2]$ nearest to
$e_1$ at which $D_2$ is selected is a limit, from the side of $e_1$, of
points whose selections give world~$2$ non-null debt.
\item \emph{Sets light in both layers} (Lemmas~\ref{lem:near}
to~\ref{lem:buys}). Near that point, (S) supplies sets in $P_2$, at
every level of the edge $[e_2,e_3]$, on which world~$2$'s full change from
credit to debt costs no more than the loss of a fixed credit.
\end{enumerate}

\emph{Outline of the proof.} Suppose $V$ is undominated. For $c>0$ let
$\UU_c$ be the set of ultrafilters on $P_2$ all of whose members are, at
world~$2$, strictly AC against credit on $Y_c\coloneqq P_3\cap
\{t^2_3<1/c\}$; Section~\ref{sec:proof} gives the precise definition.
FSE's Theorem~64 yields, for every non-null subset of a fixed non-null set,
an ultrafilter in some $\UU_c$ that contains it; so $\bigcup_c\UU_c$ is
uncountable. On the other hand, each $\UU_c$ is finite.
Theorem~\ref{thm:cut} gives a chain from $D_2$ to $D_3$ whose perturbations
are light in both layers (Section~\ref{sec:light}), and then the counting
lemma bounds $|\UU_c|$ by the length of the chain. The light sets come from
the selections near the vertex block of $e_2$ on the other edge
$[e_1,e_2]$, and Proposition~\ref{prop:rich} is what makes them available.

\emph{More worlds.} A separate note in preparation, \emph{Complete
Rankings for Any Number of Worlds}, extends Theorem~\ref{thm:main} to every
finite number $n\ge3$ of worlds. The present note proves only the case
$n=3$, and nothing in it depends on that extension.
Theorem~\ref{thm:cut} and Lemma~\ref{lem:count} hold here for every $n\ge3$;
the rest of the argument uses $n=3$ through the regions $P_x$ and through
FSE's Corollary~31 and Theorems~56 and~64, which are proved for three
worlds.

\emph{What remains open.} Incomplete rankings. TAW makes $V$ dominated,
without any completeness, when two worlds carry one-sided weak extensions of
CBOI's form and the third satisfies (IB). If some ranking is incomplete and
at most one world carries such an extension, we know neither a proof that
$V$ is dominated nor an example in which it is undominated. In particular we
do not know whether completeness at one or two of the three worlds suffices
for the conclusion of Theorem~\ref{thm:main}. Section~\ref{sec:remarks}
lists where completeness is used.

Two independent referees checked the argument of
Sections~\ref{sec:chains}, \ref{sec:count} and~\ref{sec:light} and of
Proposition~\ref{prop:finite} and Theorem~\ref{thm:main}; their corrections
are incorporated. Proposition~\ref{prop:rich} is proved here by a shorter
route than the one they checked. This proof, and the note as written, were
checked by a further independent referee, whose corrections are also
incorporated. Results that are not in the citable notes are proved in full.

\section{Setting and inputs}\label{sec:setting}

\subsection{Problems, models and axioms}\label{sec:axioms}

We use the setting and notation of FSE, Section~1, and recall what is
needed. There are $n\ge3$ worlds; from Section~\ref{sec:rich} on, $n=3$.
$\Delta\coloneqq\{p\in\R^n:p\ge0,\ \sum_ip_i=1\}$ has vertices $e_i$.
Binary problems $\{g,-g\}$ are represented by $g\in
H\coloneqq\{g\in\R^n:g_1>0\}$, with normalized Gaussian measure~$\mu$;
coordinate hyperplanes are null, and $H'$ is the set of mixed problems in
$H$. Models are measurable and identified modulo null sets. A model is
\emph{regular} if it is never contradictory and correct on unanimous
problems. Coherent models are regular; $V$ is the coherent model with credal
set $\Delta$, and $D_i$ is the ideal model of world~$i$, which recommends
the option good at~$i$ on every problem. All models below are regular.

On a mixed problem a regular model gives the option good at world~$v$
(\emph{credit} at $v$, rank~$0$), silence (rank~$1$), or the option bad at
$v$ (\emph{debt} at $v$, rank~$2$). Regular models agree on unanimous
problems, so a regular model $M$ is determined by its world-$v$ rank
function $r_v(\cdot,M)\colon H'\to\{0,1,2\}$, and every measurable function
$H'\to\{0,1,2\}$ is the world-$v$ rank function of a regular model. For
measurable $Q\subseteq H'$, the models $M+\debt_v(Q)$, $M+\cred_v(Q)$ and
$M+\sil(Q)$ agree with $M$ off $Q$ and have world-$v$ rank $2$, $0$, $1$ on
$Q$; silence does not depend on the world. $V$ is silent on $H'$, and $D_v$
has world-$v$ rank $0$ on $H'$. We write $X\le^v_\pw Y$ if
$r_v(\cdot,X)\le r_v(\cdot,Y)$ on $H'$.

At each world, $\preceq_i$ is a preorder on models; smaller is better, and
$X\prec_iY$ means that $X\preceq_iY$ and not $Y\preceq_iX$. The ranking of
world~$i$ is \emph{complete} if any two models are comparable: $X\preceq_iY$
or $Y\preceq_iX$. In a preorder,
$X\preceq_iY\prec_iZ$ and $X\prec_iY\preceq_iZ$ both imply $X\prec_iZ$.
(G) at world $i$ is guidance monotonicity: $X\le^i_\pw Y$ implies
$X\preceq_iY$ (\emph{weak (G)}), and if moreover
$r_i(\cdot,X)<r_i(\cdot,Y)$ on a non-null set, then $X\prec_iY$ (strict
(G)). (IB) at world~$i$ is independence under replacement of shared advice:
replacing advice that two models share on a measurable set of problems by
other shared advice leaves their comparison at~$i$ unchanged. We use it in
the following form.
\begin{enumerate}[label=(IB$'$),leftmargin=*]
\item If $X,Y,X',Y'$ are regular, $\Theta\coloneqq\{X\ne Y\}=\{X'\ne Y'\}$,
and $X=X'$ and $Y=Y'$ on $\Theta$, then $X\preceq_iY\iff X'\preceq_iY'$ and
$X\prec_iY\iff X'\prec_iY'$.
\end{enumerate}
Indeed, off $\Theta$ each pair shares its advice, and $(X',Y')$ arises from
$(X,Y)$ by replacing it; the strict form follows by applying the weak form
to $(X,Y)$ and to $(Y,X)$. So whether $X\preceq_vY$ depends only on the set
$\{X\ne Y\}$ and on the world-$v$ ranks of $X$ and $Y$ there.

(S) supplies coherent undominated models $C_p$, $p\in\Delta$, competitive on
$\supp p$, with closed sections $A_i(X)\coloneqq\{p:C_p\preceq_iX\}$. A model
$C$ is \emph{competitive} on a set $\mathcal W$ of worlds if for every model
$X$ some $i\in\mathcal W$ has $C\preceq_iX$. A model $M$ is \emph{dominated} if some model $X$ has
$X\preceq_iM$ at every world and $X\prec_iM$ at some world. By CBOI's
Lemma~1, a dominated model is dominated by a selected coherent model (FSE,
Section~1.1). By the face lemma (FSE, Section~1.1) the credal set of $C_p$
lies in the face spanned by $\supp p$; in particular $C_{e_i}=D_i$. For $q$
on an edge $[e_v,e_w]$, $C_q$ is a \emph{selection on the edge}; if $q$ lies
in the open edge $(e_v,e_w)$, $C_q$ is competitive on $\{v,w\}$. We use two
consequences of (S) (FSE, Section~1.1; CCF, Section~3).
\begin{enumerate}[label=(T\arabic*),leftmargin=*,itemsep=2pt]
\item If world $k$ is complete, $Y$ is regular and $Y\prec_kC_{p^*}$, then
$Y\prec_kC_p$ for all $p$ near $p^*$.
\item If $p^*\in[e_v,e_w]$, $X$ is regular and $X\prec_wC_{p^*}$, then
$C_p\preceq_vX$ for all $p$ in the open edge $(e_v,e_w)$ near $p^*$.
\end{enumerate}
(T1) holds because $A_k(Y)$ is closed and does not contain $p^*$, and
completeness turns $C_p\not\preceq_kY$ into $Y\prec_kC_p$. (T2) holds
because $A_w(X)$ is closed and does not contain $p^*$, and $C_p$ is
competitive on $\{v,w\}$. Only (T1) uses completeness.

\subsection{Debt, credit and pairs}\label{sec:pairs}

For measurable $B,C\subseteq H'$ (FSE, Section~1.2),
\begin{gather*}
 B\preceq^v_dC\ :\Longleftrightarrow\ V+\debt_v(B)\preceq_vV+\debt_v(C),\\
 B\preceq^v_cC\ :\Longleftrightarrow\ V+\cred_v(C)\preceq_vV+\cred_v(B).
\end{gather*}
So $B\preceq^v_dC$ says that debt on $B$ is no worse than debt on $C$, and
$B\preceq^v_cC$ says that credit on $B$ is worth no more than credit on $C$.
Their strict parts are written $\prec^v_d$ and $\prec^v_c$. For world~$1$ we
write $\preceq_d,\preceq_c,\prec_d,\prec_c$. Both relations are transitive.
By (IB$'$), $V$ may be replaced by any regular model silent on $B\cup C$,
and $B$ and $C$ may overlap. By weak (G), $B\subseteq C$ implies
$B\preceq^v_dC$ and $B\preceq^v_cC$.

A \emph{pair} $(A,T)$ at world $v$ consists of disjoint measurable sets
$A,T\subseteq H'$ (debt part and credit part). It is \emph{SI} (a strict
improvement) if $V+\debt_v(A)+\cred_v(T)\prec_vV$,
\emph{weakly SI} if $V+\debt_v(A)+\cred_v(T)\preceq_vV$, and \emph{strictly
AC} if $V\prec_vV+\debt_v(A)+\cred_v(T)$. If world $v$ is complete, a pair
that is not weakly SI is strictly AC. By (IB$'$), $V$ may be replaced by any
regular model silent on $A\cup T$. By strict (G), a pair $(A,T)$ with $A$
null and $T$ non-null is SI.

Let $\Alg$ be the measure algebra of $H'$. Ideals and ultrafilters are those
of $\Alg$; so ``$B$ is contained in $C$'' is read modulo null sets when
membership in an ideal is concerned. An ultrafilter \emph{avoids} an ideal if
they are disjoint. For non-null $A\subseteq H'$, $\Jid_A$ is the family of
measurable sets contained in a finite union of measurable sets each
$\preceq_dA$; it is an ideal (FSE, Definition~1 and Lemma~2).

\subsection{Edges, frames and three worlds}\label{sec:three}

For distinct worlds $v,w$ put $D_{vw}\coloneqq\{g\in H':g_vg_w<0\}$ and
$t^v_w\coloneqq|g_v|/|g_w|$ on $D_{vw}$, so $t^w_v=1/t^v_w$; the law of
$t^v_w$ under $\mu|_{D_{vw}}$ has a positive density on $(0,\infty)$, and
its level sets are null. A \emph{band} of the edge $[e_v,e_w]$ is a set
$D_{vw}\cap\{s<t^v_w<s'\}$ with $0\le s<s'\le\infty$. On $D_{vw}$ the
option good at $v$ is the option bad at $w$, so the world-$v$ and world-$w$
ranks of every regular model add up to $2$ (FSE, equation~(1)):
\begin{equation}
 \cred_v=\debt_w\quad\text{and}\quad\debt_v=\cred_w\qquad\text{on }D_{vw};
 \label{eq:swap}
\end{equation}
off $D_{vw}$ the two worlds agree, so their ranks are equal.

\emph{Edge models} (FSE, Section~1.5). By the face lemma a selection on
$[e_v,e_w]$ has a credal interval on that edge, and there are
$0\le\beta\le\alpha\le\infty$ such that it gives credit at $v$ on
$D_{vw}\cap\{t^v_w>\alpha\}$, silence on $D_{vw}\cap\{\beta<t^v_w<\alpha\}$,
debt at $v$ on $D_{vw}\cap\{t^v_w<\beta\}$, and the option good at $v$ and
$w$ off $D_{vw}$. It is written $(\alpha,\beta)$, or
$(\alpha^{(v)},\beta^{(v)})$, \emph{in world $v$'s frame}; in world $w$'s
frame it is $(\alpha^{(w)},\beta^{(w)})=(1/\beta^{(v)},1/\alpha^{(v)})$, by
\eqref{eq:swap}. Regular models of this form are \emph{edge models}. The
model $(0,0)$ is $D_v$ and $(\infty,\infty)$ is $D_w$; since $\beta\le\alpha$,
an edge model with $\alpha=0$ is $D_v$, and one with $\beta=\infty$ is $D_w$.
A selection on $[e_v,e_w]$ is \emph{two-sided} if $0<\beta\le\alpha<\infty$;
its credal set then contains neither endpoint.

\emph{Three worlds.} For $n=3$, as in FSE, Section~1.4, put
\[
 P_1\coloneqq\{g_2<0,\ g_3<0\},\qquad P_2\coloneqq\{g_2<0<g_3\},\qquad
 P_3\coloneqq\{g_3<0<g_2\}
\]
inside $H$. $P_x$ is the set of mixed problems at which world $x$ is the
odd one out, $H'=P_1\cup P_2\cup P_3$ up to null sets, and
$D_{vw}=P_v\cup P_w$. On each $P_x$ the joint law of two ratios
$|g_a|/|g_b|$ and $|g_c|/|g_d|$ with $\{a,b\}\ne\{c,d\}$ has a positive
density on $(0,\infty)^2$; so every subset of $P_x$ defined by a nonempty
open condition on one or two such ratios is non-null. For $0<b<\infty$,
$\Hi vwb\coloneqq D_{vw}\cap\{t^v_w>b\}$. In world~$1$'s frame on
$[e_1,e_2]$ we write $E_2(b)\coloneqq(\infty,b)$ for $0\le b\le\infty$: it
gives debt at $1$ on $D_{12}\cap\{t^1_2<b\}$, silence on the rest of
$D_{12}$, and the option good at $1$ and $2$ off $D_{12}$; $E_2(\infty)=D_2$.
For a selection $E$ on $[e_1,e_2]$ put $T_0(E)\coloneqq P_2\cap
\{t^1_2>\alpha^{(1)}_E\}$ (FSE, Section~4); there $E$ gives credit at $1$
and $3$, and debt at $2$.

\subsection{Results used}\label{sec:inputs}

Statements made in FSE for world~$1$ and an edge $[e_1,e_j]$ hold for every
world $v$ and edge $[e_v,e_w]$, with $D_{vw}$ and $t^v_w$ in place of
$D_{1j}$ and $t_j$, since the axioms are symmetric in the labels and every
object involved is a set of problems (FSE, Section~1.3). We use the
following results of FSE, with FSE's numbering.
\begin{enumerate}[label=(F\arabic*),leftmargin=*,itemsep=3pt]
\item\label{in:uf} (Lemma~5(a)) Every ultrafilter contains sets of
arbitrarily small measure. (Proof of Lemma~5(b)) If $\Iid$ is an ideal and
$R\notin\Iid$, some ultrafilter contains $R$ and avoids $\Iid$. The proof in
FSE is written for $\Jid_A$ but uses only that it is an ideal: the family
$\{R\}\cup\{H'\setminus J:J\in\Iid\}$ has the finite intersection property,
since $R\setminus(J_1\cup\dots\cup J_s)$ is non-null when $R\notin\Iid$, and
any ultrafilter extending it avoids $\Iid$.
\item\label{in:ends} (Lemma~6) Assume (G) at worlds $1$ and $j$ and (S),
and write $C_q=(\alpha_q,\beta_q)$ in world~$1$'s frame for
$q\in[e_1,e_j]$. (a) If $C_{q^*}=D_j$ and $q_m\to q^*$ along the edge, then
$\beta_{q_m}\to\infty$. (b) $S_j\coloneqq\{q\in[e_1,e_j]:C_q=D_j\}$ is
closed, contains $e_j$ and not $e_1$; if $q_e$ is its point nearest to
$e_1$, every $q\in[e_1,q_e)$ has $C_q\ne D_j$.
\item\label{in:yard} (Lemma~8) Let $n=3$, assume (G) at every world, (IB)
at world~$1$ and (S), let $j\in\{2,3\}$ and let $k$ be the third world.
(a) If world~$1$ is complete, then for every non-null $A\subseteq H'$ and
every open interval $I\subseteq(0,\infty)$ there is a non-null
$Z\subseteq P_1\cap\{t^1_j\in I\}$ with $Z\preceq_dA$ (namely
$Z=D_{1j}\cap D_{1k}\cap\{t^1_k>b_A\}\cap\{t^1_j\in I\}$ for a suitable
$b_A<\infty$; $D_{1j}\cap D_{1k}=P_1$). (b) For every non-null $T\subseteq H'$ there is
$a>0$ with $D_{1j}\cap\{t^1_j<a\}\preceq_cT$.
\item\label{in:edge} (Corollary~31) Let $n=3$, let all rankings be complete
with (G) and (IB), and assume (S). If $V$ is undominated, every edge carries
a two-sided selection.
\item\label{in:LH} (Theorem~56(i)) Let $n=3$, let all rankings be complete
with (G) and (IB), and assume (S). At every world $v$, for every non-null
$R\subseteq H'$ there is an SI pair $(A,T)$ of world $v$ with $A$ non-null
and $T\subseteq R$.
\item\label{in:own} (Theorem~64) Let $n=3$, let all rankings be complete
with (G) and (IB), assume (S), and suppose $V$ is undominated. Let $E$ be a
two-sided selection on $(e_1,e_2)$ and $T\subseteq T_0(E)$ non-null. Then
there is $c\in(0,\infty)$ such that $T$ is not a finite union of measurable
sets $Q$ for which $(Q,P_3\cap\Hi32c)$ is a weakly SI pair of world~$2$.
Here $P_3\cap\Hi32c=P_3\cap\{t^2_3<1/c\}$.
\end{enumerate}
We also use Theorem~64 in the relabelled form stated after its proof in
FSE. The permutation of labels $2\mapsto1$, $1\mapsto3$, $3\mapsto2$ maps
$P_2$ to $P_1$, the edge $[e_1,e_2]$ to $[e_3,e_1]$, the set $T_0(E)$, on
which $E$ gives the option good at $1$ and $3$, to the set
$P_1\cap\{t^1_3<\beta^{(1)}_E\}$, on which $E$ gives the option good at $3$
and $2$, that is, debt at~$1$, and $P_3\cap\Hi32c$ to
$P_2\cap\Hi21c=P_2\cap\{t^1_2<1/c\}$. So:
\begin{enumerate}[label=(F\arabic*$'$),leftmargin=*,start=6]
\item\label{in:ownrel} Under the hypotheses of \ref{in:own}, let $E$ be a
two-sided selection on $(e_1,e_3)$, with debt threshold $\beta^{(1)}_E$ in
world~$1$'s frame, and let $T\subseteq P_1\cap\{t^1_3<\beta^{(1)}_E\}$ be
non-null. Then there is $c\in(0,\infty)$ such that $T$ is not a finite union
of measurable sets $Q$ for which $(Q,P_2\cap\{t^1_2<1/c\})$ is a weakly SI
pair of world~$1$.
\end{enumerate}
From TAW we use only its Corollary~72, for Corollary~\ref{cor:coverage}. A
loss of \emph{CBOI's form} at world $k$ is
$L_k(M)\coloneqq\sum_{l=1}^3\nu_{kl}\{g:r_k(g,M)\ge l\}$, where the
$\nu_{kl}$ are finite nonnegative finitely additive charges on the
measurable subsets of $H$ that vanish exactly on null sets; it is a
\emph{one-sided weak extension} if
\[
 B\text{ coherent undominated},\ X\text{ regular},\ B\preceq_kX
 \ \Longrightarrow\ L_k(B)\le L_k(X).
 \tag{W$_k$}
\]
\begin{enumerate}[label=(TAW),leftmargin=*]
\item (Corollary~72) Let $n=3$, assume (G) and (S), and let every world be a
preorder satisfying (IB). If $V$ is undominated, then at most one world has
a loss of CBOI's form satisfying (W$_k$).
\end{enumerate}
TAW's Corollary~72 is its main theorem, Theorem~71, after a relabelling of
the worlds. Under (G) and (S), Theorem~71 makes $V$ dominated when worlds
$2$ and $3$ have such losses and world~$1$ satisfies (IB); it uses completeness at no world and
(IB) only at world~$1$.

\emph{Standing hypotheses} (P1): $n=3$; every ranking is a complete preorder
satisfying (G) and (IB); (S). Each result below states which of them it
uses.

\section{Connecting chains and the no-saturated-cut theorem}\label{sec:chains}

In this section and the next, $v\ne w$ are two worlds, $D\coloneqq D_{vw}$
and $t\coloneqq t^v_w$; the number $n\ge3$ of worlds is arbitrary.

\begin{definition}\label{def:chain}
Let $\Bfam$ be a family of measurable subsets of $D$. Two regular models are
\emph{$\Bfam$-neighbours} if they agree off some member of $\Bfam$. A
\emph{connecting chain} of length $\ell\ge1$ from $D_v$ to a regular model
$X$ is a sequence of regular models
\[
 D_v=Y_0,\ X_1,\ Y_1,\ X_2,\ \dots,\ Y_{\ell-1},\ X_\ell=X
\]
such that $X_i\preceq_vY_{i-1}$ for $1\le i\le\ell$ and, for
$1\le i\le\ell-1$, $Y_i$ agrees with $X_i$ off a set $B_i\in\Bfam$. The sets
$B_i$ are the \emph{perturbation sets} of the chain. The family $\Bfam$ is
\emph{band-dense} if every band of the edge $[e_v,e_w]$ contains a non-null
member of $\Bfam$.
\end{definition}

A chain alternates descents $Y_{i-1}\to X_i$ in $\preceq_v$ with
perturbations $X_i\to Y_i$ on members of $\Bfam$; a perturbation may make
the model better or worse.

\begin{lemma}\label{lem:reach}
Let $\preceq_v$ be a preorder and $\Bfam$ a family of measurable subsets of
$D$. The set $\mathcal R$ of regular models $X$ for which there is a
connecting chain from $D_v$ to $X$ contains $D_v$. It is
\emph{down-closed}: if $X\in\mathcal R$ and $Z\preceq_vX$, then
$Z\in\mathcal R$. It is \emph{$\Bfam$-saturated}: if $X\in\mathcal R$ and
$X'$ is a $\Bfam$-neighbour of $X$, then $X'\in\mathcal R$.
\end{lemma}
\begin{proof}
$Y_0=X_1=D_v$ is a chain of length $1$. If a chain ends at $X_\ell=X$ and
$Z\preceq_vX$, then $Z\preceq_vY_{\ell-1}$, and replacing $X_\ell$ by $Z$
gives a chain to $Z$. If $X'$ agrees with $X_\ell=X$ off a member of
$\Bfam$, appending $Y_\ell\coloneqq X'$ and $X_{\ell+1}\coloneqq X'$ gives a
chain to $X'$.
\end{proof}

By induction along chains, $\mathcal R$ is contained in every down-closed
$\Bfam$-saturated set of regular models that contains $D_v$. So
$D_w\in\mathcal R$ if and only if there is no \emph{saturated cut}: no
down-closed $\Bfam$-saturated set of regular models contains $D_v$ but not
$D_w$.

\begin{theorem}\label{thm:cut}
Let world $v$ be complete, and assume (G) at worlds $v$ and $w$ and (S). For
every band-dense family $\Bfam$ of measurable subsets of $D$ there is a
connecting chain from $D_v$ to $D_w$. Every connecting chain from $D_v$ to
$D_w$ has length at least $2$.
\end{theorem}
\begin{proof}
A chain of length $1$ would give $D_w=X_1\preceq_vY_0=D_v$. But $D_v$ has
world-$v$ rank $0$ on $H'$, and $D_w$ has world-$v$ rank $2$ on the non-null
set $D$; so $D_v\prec_vD_w$ by strict (G).

Suppose that there is no chain from $D_v$ to $D_w$, and let $\mathcal R$ be
as in Lemma~\ref{lem:reach}; so $D_w\notin\mathcal R$. Put
$\Sigma\coloneqq\{q\in[e_v,e_w]:C_q\in\mathcal R\}$. By the face lemma
$C_{e_v}=D_v\in\mathcal R$ and $C_{e_w}=D_w\notin\mathcal R$, so
$e_v\in\Sigma$ and $e_w\notin\Sigma$. Write $C_q=(\alpha_q,\beta_q)$ in
world $v$'s frame.

\emph{The complement of $\Sigma$ in the edge is open.} Let
$C_q\notin\mathcal R$. Then $C_q\ne D_v$, so $\alpha_q>0$, since an edge
model with $\alpha=0$ is $D_v$. The band $D\cap\{0<t<\alpha_q\}$ contains a
non-null $B\in\Bfam$, and on $B$ the model $C_q$ gives silence or debt
at~$v$. Put $M^-\coloneqq C_q+\cred_v(B)$. It is a $\Bfam$-neighbour of
$C_q$; as the neighbour relation is symmetric and $\mathcal R$ is
$\Bfam$-saturated, $M^-\notin\mathcal R$. By strict (G), $M^-\prec_vC_q$.
By (T1), which uses the completeness of world $v$, $M^-\prec_vC_{q'}$ for
all $q'$ near $q$. For such $q'$, $C_{q'}\in\mathcal R$ would give
$M^-\in\mathcal R$, since $\mathcal R$ is down-closed. So no $q'$ near $q$
lies in $\Sigma$.

\emph{$\Sigma$ is open.} Let $C_q\in\mathcal R$. Then $C_q\ne D_w$, so
$\beta_q<\infty$, since an edge model with $\beta=\infty$ is $D_w$; and
$q\ne e_w$. The band $D\cap\{\beta_q<t<\infty\}$ contains a non-null
$B\in\Bfam$, and on $B$ the model $C_q$ gives silence or credit at $v$, that
is, by \eqref{eq:swap}, silence or debt at $w$. Put
$M^+\coloneqq C_q+\debt_v(B)$, which gives credit at $w$ on $B$. Then
$M^+\in\mathcal R$ by saturation, and $M^+\prec_wC_q$ by strict (G) at
world $w$. By (T2), $C_{q'}\preceq_vM^+$ for all $q'$ in the open edge near
$q$, and these $C_{q'}$ lie in $\mathcal R$ because $\mathcal R$ is
down-closed. The only points of the edge outside the open edge are $e_v$,
which lies in $\Sigma$, and $e_w\ne q$, which we exclude by taking the
neighbourhood of $q$ small. So all points of the edge near $q$ lie in
$\Sigma$.

Thus $\Sigma$ is open and closed in the connected edge, contains $e_v$ and
not $e_w$: a contradiction.
\end{proof}

The proof uses the completeness of world $v$ only in (T1), strict (G) at
worlds $v$ and $w$, and (S) through (T1), (T2) and the face lemma. It uses
no (IB), and nothing about $V$ or about the other worlds.

\section{The counting lemma}\label{sec:count}

For a measurable $B\subseteq D$, the model $D_v+\debt_v(B)$ differs from
$D_v$ by a \emph{full change} from credit to debt at $v$ on $B$. For
disjoint $B,K\subseteq D$, by (IB$'$), the comparison of
$D_v+\debt_v(B)$ with $D_v+\debt_v(K)$ is a comparison of the full change
on $B$ with the full change on $K$.

\begin{lemma}\label{lem:count}
Let $\preceq_v$ be a preorder that satisfies weak (G) and (IB) at
world~$v$. Let
\[
 D_v=Y_0,\ X_1,\ Y_1,\ \dots,\ Y_{\ell-1},\ X_\ell=D_w
\]
be a connecting chain of length $\ell\ge2$ whose perturbation sets
$B_1,\dots,B_{\ell-1}$ are contained in a measurable set $F$. Then there
are no pairwise disjoint measurable sets $K_1,\dots,K_\ell\subseteq
D\setminus F$ such that
\begin{equation}
 D_v+\debt_v(B_i)\prec_vD_v+\debt_v(K_l)\qquad
 \text{for all }1\le i\le\ell-1\text{ and }1\le l\le\ell.
 \label{eq:heavy}
\end{equation}
\end{lemma}
\begin{proof}
Suppose such sets $K_l$ exist, and put
$\Pi_k\coloneqq D_v+\debt_v(K_1\cup\dots\cup K_k)$ for $0\le k\le\ell$. We
show by induction that $Y_k\prec_v\Pi_{k+1}$ for $0\le k\le\ell-1$.

For $k=0$, $D_v\preceq_vD_v+\debt_v(B_1)$ by weak (G), and
$D_v+\debt_v(B_1)\prec_v\Pi_1$ by \eqref{eq:heavy} with $i=l=1$; here
$\ell\ge2$ is used. So $Y_0\prec_v\Pi_1$.

Let $0\le k\le\ell-2$ and $Y_k\prec_v\Pi_{k+1}$. Then
$X_{k+1}\preceq_vY_k\prec_v\Pi_{k+1}$, so $X_{k+1}\prec_v\Pi_{k+1}$. Put
$B\coloneqq B_{k+1}$, so that $Y_{k+1}$ agrees with $X_{k+1}$ off $B$, and
\[
 X^+\coloneqq X_{k+1}+\debt_v(B),\qquad
 X^0\coloneqq X_{k+1}+\cred_v(B),\qquad
 \Pi'\coloneqq\Pi_{k+1}+\debt_v(B).
\]
Since $B\subseteq F$ and every $K_l$ misses $F$, the set $B$ misses every
$K_l$; so $\Pi_{k+1}$ has world-$v$ rank $0$ on $B$.
\begin{enumerate}[label=(\alph*),leftmargin=*,itemsep=2pt]
\item $Y_{k+1}\preceq_vX^+$ by weak (G): the two models agree off $B$, and
$X^+$ has rank $2$ on $B$.
\item $X^0\preceq_vX_{k+1}$ by weak (G), and hence $X^0\prec_v\Pi_{k+1}$.
\item $X^+\prec_v\Pi'$. Off $B$, $X^0=X^+=X_{k+1}$ and $\Pi_{k+1}=\Pi'$. On
$B$, $X^0$ and $\Pi_{k+1}$ both have rank $0$, and $X^+$ and $\Pi'$ both
have rank $2$. So the pairs $(X^0,\Pi_{k+1})$ and $(X^+,\Pi')$ have the same
difference set, which misses $B$, and agree on it; (IB$'$) carries (b)
over. The two pairs differ only in the advice they share on $B$.
\item $\Pi'\prec_v\Pi_{k+2}$. Both models have rank $2$ on
$K_1\cup\dots\cup K_{k+1}$ and rank $0$ off
$K_1\cup\dots\cup K_{k+2}\cup B$. On $B$ their ranks are $2$ and $0$, and on
$K_{k+2}$ they are $0$ and $2$. The pair
$(D_v+\debt_v(B),D_v+\debt_v(K_{k+2}))$ has the same difference set
$B\cup K_{k+2}$ and the same ranks on it. So (IB$'$) and \eqref{eq:heavy}
with $i=k+1\le\ell-1$ and $l=k+2\le\ell$ give $\Pi'\prec_v\Pi_{k+2}$.
\end{enumerate}
Hence $Y_{k+1}\preceq_vX^+\prec_v\Pi'\prec_v\Pi_{k+2}$, and the induction is
complete.

For $k=\ell-1$ we obtain $D_w=X_\ell\preceq_vY_{\ell-1}\prec_v\Pi_\ell$.
But $\Pi_\ell$ has world-$v$ rank $2$ on $K_1\cup\dots\cup K_\ell\subseteq D$
and rank $0$ elsewhere on $H'$, while $D_w$ has world-$v$ rank $2$ on $D$
and rank $0$ off $D$. So $\Pi_\ell\le^v_\pw D_w$, and $\Pi_\ell\preceq_vD_w$
by weak (G). Thus $D_w\prec_vD_w$, which is impossible.
\end{proof}

The lemma uses only transitivity, weak (G) and (IB) at world $v$: neither
completeness, nor strict (G), nor (S), nor anything about other worlds.
Under strict (G) at $v$, every connecting chain from $D_v$ to $D_w$ has
length $\ell\ge2$, by the first paragraph of the proof of
Theorem~\ref{thm:cut}.

\section{Rich far ends}\label{sec:rich}

From now on $n=3$. In this section we work in world~$1$'s frame on the edge
$[e_1,e_2]$: $t\coloneqq t^1_2$ on $D_{12}=P_1\cup P_2$, and
$C_q=(\alpha_q,\beta_q)\coloneqq(\alpha^{(1)}_q,\beta^{(1)}_q)$ for
$q\in[e_1,e_2]$. Let $q_e$ be the point of $S_2=\{q\in[e_1,e_2]:C_q=D_2\}$
nearest to $e_1$ (\ref{in:ends}(b)); $q_e\ne e_1$. We call $q_e$ the
\emph{far end of the frame $(1,2)$}: it is the end of the vertex block $S_2$
of $e_2$ that faces $e_1$. (Seen from world~$2$, whose ideal model is
selected on $S_2$, it is the near end of that block.) We parametrize the
edge by $q(s)\coloneqq(1-s)e_1+se_2$, $s\in[0,1]$, and write
$C(s)\coloneqq C_{q(s)}$, with thresholds $\alpha(s),\beta(s)$. Let
$s_e\in(0,1]$ be the parameter of $q_e$. For $0<s<s_e$, $q(s)$ lies in the
open edge and $C(s)\ne D_2$.

\begin{definition}\label{def:rich}
The far end of the frame $(1,2)$ is \emph{rich} if there are
$s_m\in(0,s_e)$ with $s_m\to s_e$ and $\alpha(s_m)<\infty$ for every $m$.
\end{definition}

In world~$2$'s frame, $\alpha(s)<\infty$ means
$\beta^{(2)}_{q(s)}=1/\alpha(s)>0$: the selection $C(s)$ gives debt at
world~$2$ on the non-null set $D_{21}\cap\{t^2_1<1/\alpha(s)\}$. If the far
end is not rich, all selections $C(s)$ with $s<s_e$ near $s_e$ have $e_2$ in
their credal sets.

\begin{proposition}\label{prop:rich}
Assume (P1) and suppose that $V$ is undominated. Then the far end of the
frame $(1,2)$ is rich.
\end{proposition}
\begin{proof}
Suppose not. Then there is $s_0\in(0,s_e)$ with $\alpha(s)=\infty$ for all
$s\in(s_0,s_e)$. For such $s$, $C(s)\ne D_2$, so $\beta(s)<\infty$, and
$C(s)=E_2(\beta(s))$. For $0\le b<b'<\infty$, the models $E_2(b)$ and
$E_2(b')$ differ exactly on the non-null band $D_{12}\cap\{b<t<b'\}$, where
$E_2(b)$ is silent and $E_2(b')$ gives debt at $1$, that is, credit at $2$
by \eqref{eq:swap}. By strict (G) at worlds $1$ and $2$,
\begin{equation}
 E_2(b)\prec_1E_2(b')\quad\text{and}\quad E_2(b')\prec_2E_2(b)
 \qquad(0\le b<b'<\infty).
 \label{eq:E2}
\end{equation}

\emph{Step 1: $\beta$ is continuous on $(s_0,s_e)$, and $\beta(s)\to\infty$
as $s\uparrow s_e$.} Let $b\in[0,\infty)$ and $s\in(s_0,s_e)$. By
\eqref{eq:E2} and reflexivity, $C(s)\preceq_1E_2(b)$ holds if and only if
$\beta(s)\le b$, and $C(s)\preceq_2E_2(b)$ holds if and only if
$\beta(s)\ge b$. The sections $A_1(E_2(b))$ and $A_2(E_2(b))$ are closed by
(S). So $\{s\in(s_0,s_e):\beta(s)\le b\}$ and
$\{s\in(s_0,s_e):\beta(s)\ge b\}$ are closed in $(s_0,s_e)$ for every
$b\in[0,\infty)$, and $\beta$ is lower and upper semicontinuous on
$(s_0,s_e)$, hence continuous. Since $C_{q_e}=D_2$, \ref{in:ends}(a) gives
$\beta(s)\to\infty$ as $s\uparrow s_e$.

\emph{Step 2: every level above some $z_0$ is crossed from both sides.}
Fix $s_1\in(s_0,s_e)$ and put $z_0\coloneqq\beta(s_1)$. Let $z>z_0$. By Step~1 there
is $s_2\in(s_1,s_e)$ with $\beta>z$ on $[s_2,s_e)$. The set
$\{s\in[s_1,s_2]:\beta(s)=z\}$ is nonempty by the intermediate value theorem
and compact by continuity; let $s^-$ and $s^+$ be its least and its greatest
element. Then $s_1<s^-\le s^+<s_2$, and, again by the intermediate value
theorem, $\beta<z$ on $[s_1,s^-)$ and $\beta>z$ on $(s^+,s_e)$. Moreover
$C(s^-)=C(s^+)=E_2(z)$.

\emph{Step 3: bands at levels above $z_0$.} Let $z>z_0$ and let $A\subseteq H'$
be non-null. We show that there is $\varepsilon>0$ with
\begin{equation}
 D_{12}\cap\{z-\varepsilon<t<z\}\preceq_dA\quad\text{and}\quad
 D_{12}\cap\{z<t<z+\varepsilon\}\preceq_dA.
 \label{eq:bands}
\end{equation}

(i) By \ref{in:yard}(a) with $j=2$ there is a non-null
$Z\subseteq D_{12}\cap\{z/4<t<z/2\}$ with $Z\preceq_dA$. Let
$X\coloneqq E_2(z)+\sil(Z)$. As $E_2(z)$ gives debt at $1$ on $Z$,
$X\prec_1E_2(z)=C(s^-)$ by strict (G), and (T1), at the complete
world~$1$, gives $X\prec_1C_q$ for all $q$ near $q(s^-)$. By continuity $\beta(s)\to z$ as $s\uparrow s^-$, so
there is $s'\in[s_1,s^-)$ with $X\prec_1C(s')$ and $b'\coloneqq\beta(s')$
in $(z/2,z)$. Put $L\coloneqq D_{12}\cap\{b'<t<z\}$ and
$X_0\coloneqq E_2(b')+\sil(Z)$. Since $Z\subseteq\{t<z/2\}\subseteq\{t<b'\}$,
we have $C(s')=E_2(b')=X_0+\debt_1(Z)$, and, up to the null level set
$\{t=b'\}$, $X=X_0+\debt_1(L)$: both give debt at $1$ exactly on
$D_{12}\cap\{t<z\}\setminus Z$, silence on the rest of $D_{12}$, and the
option good at $1$ and $2$ off $D_{12}$. So
$X_0+\debt_1(L)\prec_1X_0+\debt_1(Z)$, where $X_0$ is silent on the disjoint
sets $L$ and $Z$. By (IB$'$), $V+\debt_1(L)\prec_1V+\debt_1(Z)$, that is,
$L\prec_dZ$; so $L\preceq_dZ\preceq_dA$.

(ii) By \ref{in:yard}(a) there is a non-null
$Z'\subseteq D_{12}\cap\{z+1<t<z+2\}$ with $Z'\preceq_dA$; $E_2(z)$ is silent
on $Z'$. Let $X'\coloneqq E_2(z)+\debt_1(Z')$. By \eqref{eq:swap}, $X'$
gives credit at $2$ on $Z'$, so $X'\prec_2E_2(z)=C(s^+)$ by strict (G) at
world~$2$. By (T2) with $(v,w)=(1,2)$, $C_q\preceq_1X'$ for all $q$ in the
open edge near $q(s^+)$. By continuity $\beta(s)\to z$ as $s\downarrow s^+$,
so there is $s''\in(s^+,s_e)$ with $C(s'')\preceq_1X'$ and
$b''\coloneqq\beta(s'')$ in $(z,z+1)$. Put
$L'\coloneqq D_{12}\cap\{z<t<b''\}$. Up to a null level set,
$C(s'')=E_2(b'')=E_2(z)+\debt_1(L')$, and $E_2(z)$ is silent on the
disjoint sets $L'$ and $Z'$. So $E_2(z)+\debt_1(L')\preceq_1
E_2(z)+\debt_1(Z')$, and (IB$'$) gives $L'\preceq_dZ'\preceq_dA$.

With $\varepsilon\coloneqq\min(z-b',b''-z)$, \eqref{eq:bands} follows from
(i), (ii) and weak (G).

\emph{Step 4: a band region in every $\Jid_A$.} Put
$R\coloneqq D_{12}\cap\{z_0+1<t<z_0+2\}$ and let $A\subseteq H'$ be
non-null. For each $z\in[z_0+1,z_0+2]$, Step~3 gives $\varepsilon_z>0$ such that
$D_{12}\cap\{|t-z|<\varepsilon_z\}$ is, up to the null set $\{t=z\}$, the
union of two sets $\preceq_dA$; so it belongs to $\Jid_A$. Finitely many of
the intervals $(z-\varepsilon_z,z+\varepsilon_z)$ cover the compact interval
$[z_0+1,z_0+2]$, and $\Jid_A$ is an ideal; so $R\in\Jid_A$. This holds for every
non-null $A\subseteq H'$; in particular world~$1$'s debt order is fine on
$R$ in the sense of FSE, Definition~1.

\emph{Step 5: a contradiction.} By \ref{in:edge} there is a two-sided
selection $E$ on $(e_1,e_3)$; its debt threshold $\beta^{(1)}_E$ in
world~$1$'s frame lies in $(0,\infty)$. The set
\[
 T\coloneqq P_1\cap\{t^1_3<\beta^{(1)}_E\}\cap\{z_0+1<t^1_2<z_0+2\}
\]
is non-null, and $T\subseteq R$. By \ref{in:ownrel} there is
$c\in(0,\infty)$ such that $T$ is not a finite union of measurable sets $Q$
for which $(Q,\Gamma)$ is a weakly SI pair of world~$1$, where
$\Gamma\coloneqq P_2\cap\{t^1_2<1/c\}$. The set $\Gamma$ is non-null, so
\ref{in:LH} gives an SI pair $(A,\Gamma')$ of world~$1$ with $A$ non-null and
$\Gamma'\subseteq\Gamma$. By Step~4, $R\in\Jid_A$; so, up to a null set, $T$
is contained in a finite union $Q_1\cup\dots\cup Q_r$ of measurable sets
$Q_i\preceq_dA$. Put $Q'_i\coloneqq Q_i\cap T$, and add the null set
$T\setminus(Q'_1\cup\dots\cup Q'_r)$ to $Q'_1$, which changes no model. Then
$T=Q'_1\cup\dots\cup Q'_r$, and $Q'_i\preceq_dQ_i\preceq_dA$ by weak (G).
Each $Q'_i\subseteq P_1$ misses $\Gamma\subseteq P_2$, and $A$ misses
$\Gamma'$; so the model $V+\cred_1(\Gamma')$ is silent on $Q'_i\cup A$, and
\[
 V+\debt_1(Q'_i)+\cred_1(\Gamma)\preceq_1V+\debt_1(Q'_i)+\cred_1(\Gamma')
 \preceq_1V+\debt_1(A)+\cred_1(\Gamma')\prec_1V.
\]
The first step is weak (G) (more credit at $1$ on $\Gamma\setminus\Gamma'$),
the second is $Q'_i\preceq_dA$ in the context $V+\cred_1(\Gamma')$, by
(IB$'$), and the third says that $(A,\Gamma')$ is SI. So every pair
$(Q'_i,\Gamma)$ is SI, in particular weakly SI, and $T$ is their union. This
contradicts the choice of $c$.
\end{proof}

Steps~1--4 use (G) at every world, completeness and (IB) at world~$1$, and
(S), but nothing about $V$. They show: if every selection $C(s)$ with
$s<s_e$ near $s_e$ has $e_2$ in its credal set, then world~$1$'s debt order
is fine on a band region of $D_{12}$. Step~5 excludes this when $V$ is
undominated; it is the only step of this proof that uses \ref{in:edge},
\ref{in:LH} and \ref{in:ownrel}. Only the conclusion
$\alpha(s_m)<\infty$ is used below; whether the selections $C(s_m)$ are
two-sided plays no role.

\section{Sets light in both layers}\label{sec:light}

In this section we work at world~$2$ and on its two edges: on
$D_{21}=P_1\cup P_2$ with the ratio $t^2_1=1/t^1_2$, and on
$D_{23}=P_2\cup P_3$ with the ratio $t^2_3$. On $P_2=D_{21}\cap D_{23}$
both are defined, and their joint law has a positive density on
$(0,\infty)^2$.

\begin{definition}\label{def:light}
For measurable $W,T\subseteq H'$, the set $W$ is \emph{$T$-light} if
$W\cap T=\varnothing$ and $D_2+\debt_2(W)\preceq_2D_2+\sil(T)$.
\end{definition}

So the full change from credit to debt at world~$2$ on $W$ costs world~$2$
no more than the loss of its credit on $T$: relative to $T$, the set $W$ is
light in the credit layer and in the debt layer at once. By weak (G),
every measurable subset of a $T$-light set is $T$-light.

\begin{lemma}\label{lem:near}
Assume (P1) and suppose that $V$ is undominated. Let $c>0$ and let
$T\subseteq D_{21}\cap\{t^2_1>c\}$ be non-null. Then there is $b>0$ such
that every measurable $W\subseteq D_{21}\cap\{t^2_1<b\}$ is $T$-light.
\end{lemma}
\begin{proof}
Let $q_e$ be the far end of the frame $(1,2)$ (Section~\ref{sec:rich}). By
Proposition~\ref{prop:rich} there are $q_m$ in the open edge $(e_1,e_2)$
with $q_m\to q_e$ and $\alpha^{(1)}_{q_m}<\infty$; nothing else about these
selections is used. In world~$2$'s frame on $[e_2,e_1]$, $C_{q_m}$ has
thresholds $\alpha_m\coloneqq\alpha^{(2)}_{q_m}=1/\beta^{(1)}_{q_m}$ and
$\beta_m\coloneqq\beta^{(2)}_{q_m}=1/\alpha^{(1)}_{q_m}>0$. Since
$C_{q_e}=D_2$, \ref{in:ends}(a) gives $\beta^{(1)}_{q_m}\to\infty$, so
$\alpha_m\to0$.

Let $M_T\coloneqq D_2+\sil(T)$. On $T\subseteq D_{21}$, $D_2$ gives credit
at $2$, which by \eqref{eq:swap} is debt at $1$, while $M_T$ is silent;
elsewhere the two models agree. So $M_T\prec_1D_2=C_{q_e}$ by strict (G) at
world~$1$, and (T2) with $(v,w)=(2,1)$ gives $C_q\preceq_2M_T$ for all $q$ in
the open edge near $q_e$. Fix $m$ with $C_{q_m}\preceq_2M_T$ and
$\alpha_m<c$, and put $b\coloneqq\beta_m$; then $0<b\le\alpha_m<c$.

Let $W\subseteq D_{21}\cap\{t^2_1<b\}$ be measurable. The selection
$C_{q_m}$ gives debt at $2$ on $D_{21}\cap\{t^2_1<b\}\supseteq W$, and
$D_2+\debt_2(W)$ has world-$2$ rank $2$ on $W$ and $0$ elsewhere on $H'$.
So $D_2+\debt_2(W)\le^2_\pw C_{q_m}$, and weak (G) gives
$D_2+\debt_2(W)\preceq_2C_{q_m}\preceq_2M_T$. Finally $W\cap T=\varnothing$,
since $W\subseteq\{t^2_1<b\}$, $T\subseteq\{t^2_1>c\}$ and $b<c$.
\end{proof}

\begin{lemma}\label{lem:compose}
Assume (IB) at world~$2$. Let $W_1,\dots,W_N$ be pairwise disjoint
measurable sets, and $T_1,\dots,T_N$ pairwise disjoint measurable sets with
$T_i\cap W_k=\varnothing$ for all $i,k$. If each $W_i$ is $T_i$-light, then
$W_1\cup\dots\cup W_N$ is $(T_1\cup\dots\cup T_N)$-light.
\end{lemma}
\begin{proof}
For $0\le i\le N$ put
$M_i\coloneqq D_2+\sil(T_1\cup\dots\cup T_i)+\debt_2(W_{i+1}\cup\dots\cup
W_N)$; all the sets involved are pairwise disjoint, so $M_i$ is well
defined, $M_0=D_2+\debt_2(W_1\cup\dots\cup W_N)$ and
$M_N=D_2+\sil(T_1\cup\dots\cup T_N)$. For $0\le i<N$, the models $M_i$ and
$M_{i+1}$ differ exactly on $W_{i+1}\cup T_{i+1}$, with world-$2$ ranks $2$
and $0$ on $W_{i+1}$, and $0$ and $1$ on $T_{i+1}$. So do
$D_2+\debt_2(W_{i+1})$ and $D_2+\sil(T_{i+1})$, and the first is
$\preceq_2$ the second because $W_{i+1}$ is $T_{i+1}$-light. By (IB$'$),
$M_i\preceq_2M_{i+1}$. By transitivity $M_0\preceq_2M_N$, and
$\bigcup_iW_i$ misses $\bigcup_iT_i$.
\end{proof}

\begin{lemma}\label{lem:buys}
Assume (G) and (IB) at world~$2$. Let $W,T,Y\subseteq H'$ be pairwise
disjoint measurable sets such that $W$ is $T$-light and $T\preceq^2_cY$.
\begin{enumerate}[label=(\alph*),leftmargin=*,itemsep=2pt]
\item $(W,Y)$ is a weakly SI pair of world~$2$.
\item If moreover $K\subseteq H'$ is measurable and disjoint from
$W\cup T\cup Y$, and $(K,Y)$ is a strictly AC pair of world~$2$, then
$D_2+\debt_2(W)\prec_2D_2+\debt_2(K)$.
\end{enumerate}
\end{lemma}
\begin{proof}
(a) By lightness and weak (G),
$D_2+\debt_2(W)\preceq_2D_2+\sil(T)\preceq_2D_2+\sil(T\cup W)$. The outer
two models differ exactly on $W\cup T$, with world-$2$ ranks $2$ and $1$ on
$W$, and $0$ and $1$ on $T$; so do $V+\debt_2(W)+\cred_2(T)$ and $V$. By
(IB$'$), $V+\debt_2(W)+\cred_2(T)\preceq_2V$. By definition, $T\preceq^2_cY$
means $V+\cred_2(Y)\preceq_2V+\cred_2(T)$. The model $V+\debt_2(W)$ is silent
on $T\cup Y$, so (IB$'$) gives
$V+\debt_2(W)+\cred_2(Y)\preceq_2V+\debt_2(W)+\cred_2(T)\preceq_2V$.

(b) We show
\begin{multline*}
 D_2+\debt_2(W)\preceq_2D_2+\sil(T)\preceq_2D_2+\sil(Y)\\
 \preceq_2D_2+\sil(Y\cup K)\prec_2D_2+\debt_2(K).
\end{multline*}
The first step is lightness. In the second, the two models differ exactly on
$T\cup Y$, with ranks $1$ and $0$ on $T$, and $0$ and $1$ on $Y$; so do
$V+\cred_2(Y)$ and $V+\cred_2(T)$, and (IB$'$) turns $T\preceq^2_cY$ into
the second step. The third step is weak (G). In the fourth, the two models
differ exactly on $Y\cup K$, with ranks $1$ and $0$ on $Y$, and $1$ and $2$
on $K$; so do $V$ and $V+\debt_2(K)+\cred_2(Y)$, and (IB$'$) turns the strict
AC of $(K,Y)$ into the fourth step.
\end{proof}

Part (a) does not use $K$; it is used in Section~\ref{sec:proof} before any
such $K$ is available.

\section{Proof of the theorem}\label{sec:proof}

For $c>0$ put $Y_c\coloneqq P_3\cap\{t^2_3<1/c\}=P_3\cap\Hi32c$, a non-null
set. Let $\Iid_c$ be the family of measurable sets that are contained, up to
a null set, in a finite union of measurable sets $Q\subseteq P_2$ for which
$(Q,Y_c)$ is a weakly SI pair of world~$2$; such $Q$ miss $Y_c\subseteq
P_3$. $\Iid_c$ is an ideal: it contains the null sets and is closed under
measurable subsets and finite unions. If $c\le c'$, then $Y_{c'}\subseteq
Y_c$, and a weakly SI pair $(Q,Y_{c'})$ gives a weakly SI pair $(Q,Y_c)$ by
weak (G) at world~$2$ (more credit); so $\Iid_{c'}\subseteq\Iid_c$. Let
$\UU_c$ be the set of ultrafilters that contain $P_2$ and avoid $\Iid_c$.

If world~$2$ is complete, $\Uf\in\UU_c$, $K\in\Uf$ and $K\subseteq P_2$,
then $(K,Y_c)$ is a strictly AC pair of world~$2$. Indeed $K\notin\Iid_c$,
so $(K,Y_c)$ is not weakly SI, since otherwise $K$ would cover itself; and
by completeness a pair that is not weakly SI is strictly AC.

\begin{proposition}\label{prop:finite}
Assume (P1) and suppose that $V$ is undominated. Then $\UU_c$ is finite for
every $c>0$.
\end{proposition}
\begin{proof}
Fix $c>0$.

\emph{Credit witnesses.} By \ref{in:yard}(b), applied at world~$2$ with
$j=1$ to the non-null set $Y_c$, there is $a>0$ with
$D_{21}\cap\{t^2_1<a\}\preceq^2_cY_c$. Put
$T\coloneqq P_1\cap\{a/2<t^2_1<a\}$, a non-null set. By weak (G) and
transitivity, $T'\preceq^2_cY_c$ for every measurable $T'\subseteq T$.
Split $T$ into pairwise disjoint non-null measurable sets $T_k$, $k\ge1$;
this is possible since $\mu$ is atomless. Each $T_k$ lies in
$D_{21}\cap\{t^2_1>a/2\}$, so Lemma~\ref{lem:near} gives $b_k>0$ such that
every measurable subset of $D_{21}\cap\{t^2_1<b_k\}$ is $T_k$-light.

\emph{A band-dense family of light sets.} Choose reals $\varepsilon_k$ with
$0<\varepsilon_1<b_1$ and $0<\varepsilon_k<\min(b_k,\varepsilon_{k-1}/2)$
for $k\ge2$. Let $(r_k,r'_k)_{k\ge1}$ enumerate the pairs of rationals with
$0<r_k<r'_k$, and put
\[
 W_k\coloneqq P_2\cap\{r_k<t^2_3<r'_k\}\cap
 \{\varepsilon_k/2<t^2_1<\varepsilon_k\}.
\]
Each $W_k$ is non-null, being defined by a nonempty open condition on
$(t^2_1,t^2_3)$ in $P_2$. The $W_k$ are pairwise disjoint, since their
ranges of $t^2_1$ are. Each $W_k\subseteq P_2$ misses $T\subseteq P_1$ and
$Y_c\subseteq P_3$, and $W_k$ is $T_k$-light, since
$W_k\subseteq D_{21}\cap\{t^2_1<b_k\}$. Let $\Bfam$ be the family of
measurable subsets of the sets $W_k$. It is band-dense for the edge
$[e_2,e_3]$ in world~$2$'s frame: a band $D_{23}\cap\{s<t^2_3<s'\}$, with
$0\le s<s'\le\infty$, contains every $W_k$ with $s\le r_k<r'_k\le s'$, and
there are such $k$.

\emph{The chain.} By Theorem~\ref{thm:cut} with $(v,w)=(2,3)$ (world~$2$ is
complete, (G) holds at worlds $2$ and $3$, and (S) holds), there is a
connecting chain from $D_2$ to $D_3$ with perturbation sets in $\Bfam$. Let
$\ell\ge2$ be its length and $B_1,\dots,B_{\ell-1}$ its perturbation sets.
Each $B_i$ lies in some $W_{k_i}$; let $N\coloneqq\max_ik_i$ and
\[
 F\coloneqq W_1\cup\dots\cup W_N,\qquad T^N\coloneqq T_1\cup\dots\cup T_N.
\]
By Lemma~\ref{lem:compose}, $F$ is $T^N$-light, and so is every
$B_i\subseteq F$. Moreover $T^N\preceq^2_cY_c$, and $F\subseteq P_2$,
$T^N\subseteq P_1$ and $Y_c\subseteq P_3$ are pairwise disjoint. By
Lemma~\ref{lem:buys}(a), $(F,Y_c)$ is a weakly SI pair of world~$2$; as
$F\subseteq P_2$, $F\in\Iid_c$.

\emph{Counting.} Suppose that $\UU_c$ has $\ell$ distinct members
$\Uf^1,\dots,\Uf^\ell$. Each contains $P_2$ and not $F$, hence contains
$P_2\setminus F$. For $l<l'$ choose $G_{ll'}\in\Uf^l$ with
$H'\setminus G_{ll'}\in\Uf^{l'}$, and put $G_{l'l}\coloneqq H'\setminus
G_{ll'}$; then $G_{ll'}\in\Uf^l$ for all $l\ne l'$. The sets
\[
 K_l\coloneqq(P_2\setminus F)\cap\bigcap_{l'\ne l}G_{ll'}\in\Uf^l
\]
are pairwise disjoint: for $l\ne l'$, $K_l\subseteq G_{ll'}$ and
$K_{l'}\subseteq G_{l'l}=H'\setminus G_{ll'}$. Each $K_l$ is a subset of
$P_2$ in $\Uf^l\in\UU_c$, so $(K_l,Y_c)$ is a strictly AC pair of world~$2$.
For all $i$ and $l$, the sets $B_i\subseteq F$, $T^N\subseteq P_1$,
$Y_c\subseteq P_3$ and $K_l\subseteq P_2\setminus F$ are pairwise disjoint,
$B_i$ is $T^N$-light, and $T^N\preceq^2_cY_c$. So Lemma~\ref{lem:buys}(b)
gives
\[
 D_2+\debt_2(B_i)\prec_2D_2+\debt_2(K_l)\qquad
 \text{for all }1\le i\le\ell-1\text{ and }1\le l\le\ell.
\]
Since $K_l\subseteq P_2\setminus F\subseteq D_{23}\setminus F$, this
contradicts Lemma~\ref{lem:count} with $(v,w)=(2,3)$. Hence
$|\UU_c|\le\ell-1$.
\end{proof}

\begin{theorem}\label{thm:main}
Let $n=3$, let every ranking be a complete preorder satisfying (G) and
(IB), and assume (S). Then $V$ is dominated, and hence dominated by a
selected coherent model.
\end{theorem}
\begin{proof}
Suppose that $V$ is undominated. By Proposition~\ref{prop:finite}, the set
$\UU\coloneqq\bigcup_{c\in\mathbb N,\,c\ge1}\UU_c$ is countable; list its
members as $\Uf^1,\Uf^2,\dots$ (the list may be finite or empty). By
\ref{in:edge} there is a two-sided selection $E$ on $(e_1,e_2)$. As
$\alpha^{(1)}_E<\infty$, the set $T_0(E)=P_2\cap\{t^1_2>\alpha^{(1)}_E\}$ is
non-null. By \ref{in:uf} choose $B^i\in\Uf^i$ with
$\mu(B^i)<2^{-i-1}\mu(T_0(E))$, and put
$T^*\coloneqq T_0(E)\setminus\bigcup_iB^i$. Then
$\mu(T^*)\ge\mu(T_0(E))/2>0$, and $T^*\cap B^i=\varnothing$ for every $i$.

By \ref{in:own} there is $c_0\in(0,\infty)$ such that $T^*$ is not a finite
union of measurable sets $Q$ for which $(Q,Y_{c_0})$ is weakly SI at
world~$2$. Let $c\ge c_0$ be an integer. Then $T^*$ is not a finite union of
measurable sets $Q$ for which $(Q,Y_c)$ is weakly SI. Indeed, such $Q$
satisfy $Q\subseteq T^*\subseteq P_2$, and since $Y_c\subseteq Y_{c_0}$,
weak (G) at world~$2$ (more credit) would give
$V+\debt_2(Q)+\cred_2(Y_{c_0})\preceq_2V+\debt_2(Q)+\cred_2(Y_c)\preceq_2V$:
each $(Q,Y_{c_0})$ would be weakly SI.

Next, $T^*\notin\Iid_c$. Otherwise, up to a null set,
$T^*\subseteq Q_1\cup\dots\cup Q_r$ with $Q_i\subseteq P_2$ and $(Q_i,Y_c)$
weakly SI; here $r\ge1$, since $T^*$ is non-null. Then $T^*$ is the union of
the sets $Q_i\cap T^*$ together with a null set, which we add to
$Q_1\cap T^*$; this changes no model. Each $(Q_i\cap T^*,Y_c)$ is weakly SI
by weak (G) at world~$2$ (less debt). This contradicts the previous
paragraph.

By \ref{in:uf} there is an ultrafilter $\Uf$ that contains $T^*$ and avoids
$\Iid_c$. It contains $P_2\supseteq T^*$, so $\Uf\in\UU_c\subseteq\UU$, and
$\Uf=\Uf^i$ for some $i$. Then $\Uf$ contains the disjoint sets $T^*$ and
$B^i$, which is impossible. So $V$ is dominated, and by CBOI's Lemma~1
(FSE, Section~1.1) it is dominated by a selected coherent model.
\end{proof}

\begin{corollary}\label{cor:coverage}
Let $n=3$, assume (G) and (IB) at every world, and assume (S). Suppose that
$V$ is undominated; this holds in particular if every coherent model is
undominated. Then:
\begin{enumerate}[label=(\alph*),leftmargin=*,itemsep=2pt]
\item the ranking of some world is incomplete: that world has two models
neither of which it ranks at least as good as the other;
\item at most one world has a loss of CBOI's form satisfying the one-sided
bound (W$_k$).
\end{enumerate}
\end{corollary}
\begin{proof}
(a) is Theorem~\ref{thm:main}: if all three rankings were complete, $V$
would be dominated. (b) is TAW's Corollary~72.
\end{proof}

So with three worlds, rankings that satisfy (G), (IB) and (S) and leave
every coherent model undominated must contain incomparabilities, and at
least two of the three worlds lack one-sided weak extensions of CBOI's
form.

\section{Remarks}\label{sec:remarks}

\begin{remark}\label{rem:hyp}
\emph{Where the hypotheses are used.}
\begin{enumerate}[label=(\alph*),leftmargin=*,itemsep=2pt]
\item \emph{Completeness.} At world~$2$: in Theorem~\ref{thm:cut} with
$(v,w)=(2,3)$, through (T1), and in reading the members of the
ultrafilters in $\UU_c$ as strictly AC against $Y_c$. At world~$1$: in
Proposition~\ref{prop:rich}, through (T1) and \ref{in:yard}(a). At every
world: through \ref{in:edge}, \ref{in:LH} and \ref{in:own}, which FSE proves
for complete rankings. The completeness of world~$3$ enters only through
these results of FSE. Lemmas~\ref{lem:count}, \ref{lem:compose}
and~\ref{lem:buys}, \ref{in:yard}(b) and (T2) do not use completeness, and
neither does the proof of Lemma~\ref{lem:near} once
Proposition~\ref{prop:rich} is given.
\item \emph{(IB).} At world~$2$: in Lemma~\ref{lem:count} (with $v=2$),
Lemmas~\ref{lem:compose} and~\ref{lem:buys}, and \ref{in:yard}(b). At
world~$1$: in Steps~3 and~5 of the proof of Proposition~\ref{prop:rich} and
in \ref{in:yard}(a). At every world: through \ref{in:edge}, \ref{in:LH} and
\ref{in:own}. Theorem~\ref{thm:cut} and the proof of Lemma~\ref{lem:near}
use no (IB).
\item \emph{(G).} Strict (G) is used at worlds $2$ and $3$ in
Theorem~\ref{thm:cut}, at worlds $1$ and $2$ in
Proposition~\ref{prop:rich}, at world~$1$ in Lemma~\ref{lem:near}, and
inside the results of FSE. Lemmas~\ref{lem:count}, \ref{lem:compose}
and~\ref{lem:buys} need only weak (G).
\item \emph{$V$ undominated.} This is the reductio hypothesis of
Theorem~\ref{thm:main}. It enters only through \ref{in:edge}, \ref{in:own}
and \ref{in:ownrel}: in Step~5 of the proof of Proposition~\ref{prop:rich},
hence in Lemma~\ref{lem:near} and Proposition~\ref{prop:finite}, and in the
proof of Theorem~\ref{thm:main}. Theorem~\ref{thm:cut},
Lemmas~\ref{lem:count}, \ref{lem:compose} and~\ref{lem:buys}, and Steps~1--4
of the proof of Proposition~\ref{prop:rich} assume nothing about $V$.
\item \emph{Three worlds.} The proof uses $n=3$ through the regions $P_x$
and through \ref{in:yard}--\ref{in:own}. In Section~\ref{sec:proof} the
regions keep the perturbation sets and heavy sets ($P_2$), the credit
witnesses $T_k$ ($P_1$) and the credit set $Y_c$ ($P_3$) disjoint; in Step~5
of the proof of Proposition~\ref{prop:rich} they keep the set $T$ and the
pieces $Q'_i$ ($P_1$) disjoint from the credit sets $\Gamma'\subseteq\Gamma$
($P_2$).
\end{enumerate}
\end{remark}

\begin{remark}\label{rem:count}
\emph{The counting lemma.} Lemma~\ref{lem:count} concerns a single preorder
and a single chain, and needs only transitivity, weak (G) and (IB). It does
not bound the change from $D_v$ to $D_w$ by the sum of the perturbations;
perturbations are not shared advice, so (IB) cannot add them up. Instead, at
each step it lifts both sides of the current strict comparison to common
debt on the current perturbation set, which is a change of shared advice
(step~(c)), and then trades this debt for one fresh heavy set (step~(d)).
Each perturbation costs exactly one heavy set. The price is that
\eqref{eq:heavy} compares \emph{full} changes from credit to debt; this is
why the perturbations of Section~\ref{sec:proof} must be light in both
layers at once, which Lemma~\ref{lem:near} provides. For a world ranked
exactly by a loss $\Lambda(M)=\nu_1\{r_v(\cdot,M)\ge1\}+
\nu_2\{r_v(\cdot,M)\ge2\}$, with finite finitely additive charges on $H'$
that vanish exactly on null sets (so that $\Lambda$ is defined on models
modulo null sets), the lemma reduces to a comparison of masses. With
$\varphi\coloneqq\nu_1+\nu_2$, a descent does not raise $\Lambda$ and the
$i$-th perturbation raises it by at most $\varphi(B_i)$; so
$\varphi(D)=\Lambda(D_w)-\Lambda(D_v)\le(\ell-1)\max_i\varphi(B_i)$, whereas
$\ell$ pairwise disjoint sets $K_l\subseteq D$ as in \eqref{eq:heavy} would
give $\varphi(D)>\ell\max_i\varphi(B_i)$.
\end{remark}

\begin{remark}\label{rem:S}
\emph{What (S) contributes.} The finiteness of $\UU_c$ rests on (S) in
three ways: through the chain of Theorem~\ref{thm:cut} on the edge
$[e_2,e_3]$; through the light sets of Lemma~\ref{lem:near}, which come from
the selections near the vertex block of $e_2$ on the other edge $[e_1,e_2]$
and are available because that block's far end is rich
(Proposition~\ref{prop:rich}); and through the results of FSE. World~$2$'s
credit order is light at low levels of $t^2_1$ already by \ref{in:yard}(b),
which is how the witnesses $T_k$ are found. What the selections near the
vertex block add is lightness in the debt layer at the same time, at every
level of $t^2_3$.
\end{remark}

\end{document}
